\subsection{Excluding closed systems inside prefixes}
\label{sec:projective-factor-bound}

The forbidden pair in the preceding estimate is the smallest member
of a family of forbidden factors. Keeping the whole family gives
deterministic time and space \(O(n^{O(1)}\,1.7271^n)\) for the same
transfer.

Call a nonempty factor of a paired sign word a Dyck pair when its two
coordinate words are both Dyck words. Whatever the state at the
start of the factor, every closing in either coordinate then consumes
a local opening made inside the factor. Hence every site in this
interval has both its matching edge and its antipodal edge inside
the interval. At least one component has closed by its end.
Consequently a prefix stored for further transfer contains no Dyck
pair. A terminal move that closes the entire diagram contributes
directly to the answer and creates no stored successor.

Let \(I_j\) count the paired Dyck words of length \(2j\) containing no
proper nonempty Dyck pair. These are exactly the irreducible
meandric systems of Lando and Zvonkin. A contiguous interval carries
an independent system precisely when both coordinate matchings
restrict to perfect matchings of that interval. In the two sign
words this is equivalent to a Dyck pair. Put
\[
 B(x)=\sum_{j\ge0}C_j^2x^j,\qquad I(x)=1+\sum_{j\ge1} I_j x^j.
\]
Lando and Zvonkin proved \(B(x)=I(xB(x)^2)\) \cite{lz}.

Two minimal forbidden factors cannot overlap properly. To see this,
use boundary-time intervals \([a,c]\) and \([b,d]\), with
\(a<b<c<d\), and let \(v(t)\) be the vector of the two heights.
The Dyck condition gives \(v(b)\ge v(a)=v(c)\ge v(b)\)
coordinatewise. The interval \([b,c]\) is therefore a proper
nonempty Dyck pair, a contradiction. Containment of one minimal
factor in another is excluded by definition.

Choose \(u>1\), give a letter \((i,j)\) weight \(u^{i+j}\), and set
\(S=(u+u^{-1})^2\). On any one interval of fixed coordinate sign
changes in the minimum-reflection argument, each forbidden word is
transformed by the same alphabet permutation. Nonoverlap is
preserved, and the weight of each transformed forbidden word is
one because both coordinate sums are zero. Let \(F_t\) be the
total weight of words of length \(t\) avoiding these factors.
Appending one letter to a valid word either leaves it valid or
completes its unique minimal forbidden suffix. Conversely, appending
a minimal forbidden word to a valid prefix creates the first
occurrence at the final letter, by nonoverlap. Thus
\[
 \sum_{t\ge0}F_t z^t=\frac1{I(z^2)-S z}.
\]

The first positive pole is at \(R=B(S^{-2})/S\). Indeed,
\(v\mapsto vB(v^2)\) increases strictly for \(0\le v<1/4\),
and substitution \(z=vB(v^2)\) gives denominator
\(B(v^2)(1-Sv)\). Since \(1/S<1/4\), the series \(I(z^2)\)
converges beyond \(R\) by the Lando--Zvonkin identity and
nonnegativity. The denominator is positive on \([0,R)\) and has a
simple zero at \(R\). The generating function has nonnegative
coefficients, so Pringsheim's theorem gives radius \(R\). Evaluating
at \(R(1-1/t)\) for \(t\ge2\) then gives \(F_t=O(tR^{-t})\).

Write \(\lambda=1/R\). Multiplication of the estimates for the at
most three sign intervals bounds the number of histories for fixed
parameters by \(n^{O(1)}\lambda^t u^{-D}\). The independent
bit-length estimate is \(4\,2^D\). Balancing the two bounds at
\((2u)^D=\lambda^t\) gives exponent \(\alpha t\) in base two,
where \(\alpha=\log(\lambda)/\log(2u)\). Comparison with the
remaining-length exponent \(2(n-t)\) gives the base
\(2^{2\alpha/(2+\alpha)}\).

For an exact numerical bound take \(u=41/20\),
\(S=4330561/672400\), and
\(B_5=\sum_{j=0}^5 C_j^2S^{-2j}<B(S^{-2})\).
The two rational comparisons
\[
 (S/B_5)^{156}<(41/10)^{203},\qquad
 2^{406}10000^{515}<17271^{515}
\]
give \(\alpha<203/156\) and the stated upper bound \(1.7271\).
